\ifdefined\gkver\else\def\gkver{student}\fi
\documentclass[\gkver]{gaokaozhenti}
\begin{document}

\chapter{2026年江西}

\begin{enumerate}
\item
%% number: 1
%% paperName: 2026年江西高考第 1 题 4 分
%% typeId: 1
%% type: 单选题
%% chapter: 原子和原子核
%% point: 人工核反应
%% method: 无
%% score: 4
%% degree: 950
%% duplicateId: 0
%% body:
东华加速器中子源科学装置通过 $\gamma$ 射线与钨核反应产生中子，其方程为 \ce{\gamma{} + ^{184}_{74}W -> ^{A}_{Z}X + ^{1}_{0}n}，其中 $\gamma$ 光子的静质量为零。关于方程中 $Z$ 与 $A$ 的值，下列选项正确的是\xzanswer{C}
\fourchoices[answer=C]
{$Z=73$，$A=183$}
{$Z=73$，$A=184$}
{$Z=74$，$A=183$}
{$Z=74$，$A=184$}
%% answer: C
%% memo:
\memoanswer{
核反应过程遵循电荷数守恒和质量数守恒：根据电荷数守恒，反应前总电荷数为钨核的电荷数 74，$\gamma$ 光子不带电，反应后中子电荷数为 0，有 $74=Z+0$，得 $Z=74$。

质量数守恒：反应前总质量数为钨核的质量数 184，$\gamma$ 光子静质量为 0，反应后中子质量数为 1，因此 $184=A+1$，得 $A=183$。

故选 C。
}

\item
%% number: 2
%% paperName: 2026年江西高考第 2 题 4 分
%% typeId: 1
%% type: 单选题
%% chapter: 运动和力
%% point: 动量定理
%% method: 无
%% score: 4
%% degree: 850
%% duplicateId: 0
%% body:
在篮球运动中，运动员跳起投篮后竖直落地。设地面对运动员的平均冲击力大小为 $F$，冲击力作用的时间为 $\Delta t$。运动员在触地时往往会主动屈膝下蹲，关于这一动作对 $\Delta t$ 和 $F$ 产生的影响，下列选项正确的是\xzanswer{A}
\fourchoices[answer=A]
{$\Delta t$ 增加，$F$ 减小}
{$\Delta t$ 增加，$F$ 增加}
{$\Delta t$ 减小，$F$ 减小}
{$\Delta t$ 减小，$F$ 增加}
%% answer: A
%% memo:
\memoanswer{
运动员竖直落地过程，接触地面瞬间的速度大小为 $v$，末速度为 0，整个过程动量变化量大小为 $\Delta p=mv$，为定值。取向上为正方向，根据动量定理
\[
(F-mg)\Delta t=0-(-mv)=mv
\]
整理得平均冲击力
\[
F=mg+\frac{mv}{\Delta t}
\]
运动员屈膝下蹲时，会延长速度从 $v$ 减为 0 的作用时间 $\Delta t$。由上式可知，$\Delta t$ 增大时，$\frac{mv}{\Delta t}$ 减小，因此平均冲击力 $F$ 减小。

故选 A。
}

\item
%% number: 3
%% paperName: 2026年江西高考第 3 题 4 分
%% typeId: 1
%% type: 单选题
%% chapter: 电场
%% point: 等势面
%% method: 无
%% score: 4
%% degree: 550
%% duplicateId: 0
%% body:
两块可视为无限大非平行金属板 M、N 接在恒定电压的电源上，其截面如图所示，虚线表示等势线，相邻等势线间的电势差相等。若一电子仅受电场力作用从 $P$ 处运动到 $Q$ 处，则关于电场线分布和电子运动轨迹，下列选项可能正确的是\xzanswer{B}
\onepicture[w=5.2cm]{\includesvg{figs/03a}}
\fourchoices[answer=B, ispicture=true, h=3.1cm]
{\includesvg{figs/03b}}
{\includesvg{figs/03c}}
{\includesvg{figs/03d}}
{\includesvg{figs/03e}}
%% answer: B
%% memo:
\memoanswer{
电场线分布满足的条件：电场线始终与等势线垂直，且由高电势的正极板 M 指向低电势的负极板 N。两块非平行金属板接恒定电压，板间距从左到右逐渐增大，根据 $E=\frac{U}{d}$ 定性分析可知从左到右电场变弱，则电场线的分布应该左密右疏。

电子运动轨迹需要满足：电子带负电，所受电场力方向与电场强度方向相反，即电场力指向电势升高的方向（指向正极板 M 一侧）；同时曲线运动的轨迹会向合力（电场力）的一侧弯曲。

综上可知 B 正确，ACD 错误。

故选 B。
}

\item
%% number: 4
%% paperName: 2026年江西高考第 4 题 4 分
%% typeId: 1
%% type: 单选题
%% chapter: 曲线运动
%% point: 人造地球卫星
%% method: 无
%% score: 4
%% degree: 650
%% duplicateId: 0
%% body:
我国长征八号运载火箭以“一箭十八星”的方式将千帆星座组网卫星成功送入预定轨道。已知地球质量 $M=6.0\times10^{24}\,\Ukg$，地球半径 $R=6.4\times10^{6}\,\Um$，引力常量 $G=6.67\times10^{-11}\,\UNmqkgq$，不考虑地球自转的影响，则某颗卫星在高度约为 $h=600\,\Ukm$ 的绕地圆形轨道上运行时的速度约为\xzanswer{C}
\fourchoices[answer=C]
{$6.5\,\Ukms$}
{$7.0\,\Ukms$}
{$7.6\,\Ukms$}
{$8.1\,\Ukms$}
%% answer: C
%% memo:
\memoanswer{
卫星绕地球做匀速圆周运动，万有引力提供向心力，由
\[
G\frac{Mm}{(R+h)^2}=m\frac{v^2}{R+h}
\]
可得线速度表达式
\[
v=\sqrt{\frac{GM}{R+h}}
\]
代入数据计算：轨道半径
\[
R+h=6.4\times10^{6}\,\Um+6\times10^{5}\,\Um=7.0\times10^{6}\,\Um
\]
\[
GM=6.67\times10^{-11}\times6.0\times10^{24}\approx4.0\times10^{14}\,\UNmqkg
\]
代入可得
\[
v=\sqrt{\frac{4.0\times10^{14}}{7.0\times10^{6}}}\,\Ums\approx7.6\times10^{3}\,\Ums=7.6\,\Ukms
\]
故选 C。
}

\item
%% number: 5
%% paperName: 2026年江西高考第 5 题 4 分
%% typeId: 1
%% type: 单选题
%% chapter: 直线运动
%% point: 运动图象
%% method: 无
%% score: 4
%% degree: 650
%% duplicateId: 0
%% body:
具身智能机器人应用到精密制造生产线时，需要避免其加速度突变，以显著降低机械冲击。为实现这一目的，关于机器人 $v-t$ 图像的设计，下列选项最符合要求的是\xzanswer{B}
\fourchoices[answer=B, ispicture=true, h=2.5cm]
{\includesvg{figs/05a}}
{\includesvg{figs/05b}}
{\includesvg{figs/05c}}
{\includesvg{figs/05d}}
%% answer: B
%% memo:
\memoanswer{
$v-t$ 图像的斜率代表加速度。要求避免加速度突变，也就是要求 $v-t$ 图像的斜率不能出现突然的跳变，全程加速度要连续平缓变化。

A．加速、减速阶段都是直线，斜率恒定，在速度从 0 开始加速、以及从匀速阶段切入减速的拐点处，斜率会发生突变，对应加速度突变，不符合要求，故 A 错误；

B．全程 $v-t$ 图像是光滑的连续曲线，斜率（加速度）始终是连续平缓变化的，不存在加速度的突然跳变，完全符合避免加速度突变的要求，故 B 正确；

C．加速段的后半段是倾斜直线、减速段的前半段也是倾斜直线，在直线和曲线衔接的位置，斜率会发生突变，对应加速度突变，不符合要求，故 C 错误；

D．加速段是上凸的曲线，斜率（加速度）在加速过程持续减小到 0，减速段是下凸的曲线，斜率（加速度的绝对值）持续增大再减小，在加速和匀速的衔接处、匀速和减速的衔接处，斜率会出现突变，对应加速度突变，不符合要求，故 D 错误。

故选 B。
}

\item
%% number: 6
%% paperName: 2026年江西高考第 6 题 4 分
%% typeId: 1
%% type: 单选题
%% chapter: 功和能
%% point: 动能势能
%% method: 建立模型
%% score: 4
%% degree: 450
%% duplicateId: 0
%% body:
某地需设计一额定功率为 $3.14\times10^{6}\,\UW$ 的水平轴风力发电机。在风速为 $13\Ums$、空气的密度为 $1.3\Ukgmc$、发电机叶片的旋转平面与风向垂直，以及叶片转动形成圆面内的风能有 $20\%$ 转化为电能的条件下，发电机以额定功率工作。若叶片转速为 $15\Urmin$，则该圆面的直径和叶片尖端的速度分别约为\xzanswer{D}
\fourchoices[answer=D]
{$59\,\Um$，$16\,\Ums$}
{$59\,\Um$，$46\,\Ums$}
{$118\,\Um$，$46\,\Ums$}
{$118\,\Um$，$93\,\Ums$}
%% answer: D
%% memo:
\memoanswer{
取 $t$ 时间内垂直吹过叶片旋转圆面的空气为研究对象，这部分空气的体积为
\[
V=Svt
\]
其中 $S$ 是叶片旋转圆面积，$v$ 为风速，对应空气质量
\[
m=\rho V=\rho Svt
\]
这部分空气的总动能为
\[
E_k=\frac{1}{2}mv^2=\frac{1}{2}\rho Sv^3t
\]
结合 $20\%$ 的风能转化为电能，额定功率 $P$ 满足
\[
P=\frac{\eta E_k}{t}=\frac{1}{2}\eta\rho Sv^3
\]
圆面积公式
\[
S=\pi\left(\frac{D}{2}\right)^2
\]
代入数据解得圆面的直径
\[
D\approx118\,\Um
\]
叶片转速
\[
n=15\,\Urmin=0.25\,\Urs
\]
叶片尖端做圆周运动的线速度满足
\[
v_{\text{尖}}=\omega\cdot\frac{D}{2}=2\pi n\cdot\frac{D}{2}=\pi nD
\]
代入数据可得
\[
v_{\text{尖}}\approx93\,\Ums
\]
故选 D。
}

\item
%% number: 7
%% paperName: 2026年江西高考第 7 题 4 分
%% typeId: 1
%% type: 单选题
%% chapter: 力物体的平衡
%% point: 多力平衡
%% method: 无
%% score: 4
%% degree: 450
%% duplicateId: 0
%% body:
某大桥是一座独柱斜塔斜拉桥，如图（a）所示。大桥斜塔、正索区及两组对称分布的背索区结构分别等效为刚性杆 $PO$、拉索 $PA$、$PB$ 和 $PC$，如图（b）所示，$AD$ 为桥面中心线，$BC$ 与 $AD$ 垂直，$PO$、$PA$ 所在平面与桥面垂直，且 $PO$、$PA$ 与桥面的夹角分别为 $\theta$、$\alpha$，$PB$ 与 $PC$ 长度相等且夹角为 $2\beta$，$\angle DPO=\gamma$。设 $PA$、$PB$ 和 $PC$ 对 $PO$ 的拉力分别为 $T_1$、$T_2$ 和 $T_3$。若此三力的合力沿 $PO$ 方向，且 $T_2=T_3$，不考虑拉索的质量，则 $T_1$ 与 $T_2$ 的比值为\xzanswer{A}
\onepicture[w=10cm]{figs/07.png}
\fourchoices[answer=A]
{$\frac{2\sin\gamma\cos\beta}{\sin(\theta-\alpha)}$}
{$\frac{2\sin(\theta-\alpha)\cos\beta}{\sin\gamma}$}
{$\frac{2\sin\gamma}{\sin(\theta-\alpha)\cos\beta}$}
{$\frac{2\sin(\theta-\alpha)}{\sin\gamma\cos\beta}$}
%% answer: A
%% memo:
\memoanswer{
因为 $T_2=T_3$，且 $PB$、$PC$ 关于 $POD$ 平面对称，二者的合力会沿 $PD$ 方向，大小为
\[
F_{23}=2T_2\cos\beta
\]
合力沿 $PO$ 方向，因此垂直于 $PO$ 方向的分力总和为 0，根据几何关系可得正索拉力 $T_1$ 在垂直 $PO$ 方向的分力大小为 $T_1\sin(\theta-\alpha)$。

背索合力 $F_{23}$ 在垂直 $PO$ 方向的分力大小为
\[
F_{23}\sin\gamma=2T_2\cos\beta\sin\gamma
\]
两个垂直分力大小相等、方向相反，因此
\[
T_1\sin(\theta-\alpha)=2T_2\sin\gamma\cos\beta
\]
可得
\[
\frac{T_1}{T_2}=\frac{2\sin\gamma\cos\beta}{\sin(\theta-\alpha)}
\]
故选 A。
}

\item
%% number: 8
%% paperName: 2026年江西高考第 8 题 6 分
%% typeId: 2
%% type: 多选题
%% chapter: 电磁感应
%% point: 自感与互感，涡流
%% method: 无
%% score: 6
%% degree: 650
%% duplicateId: 0
%% body:
某小组探究线圈的自感现象，实验电路如图所示，其中 $A_1$ 和 $A_2$ 是两个相同的灯泡，线圈的自感系数很大。下列选项正确的是\xzanswer{BD}
\onepicture[w=6cm]{figs/08.png}
\fourchoices[answer=BD]
{闭合 S，稳定后缓慢减小 $R$ 的阻值，$A_1$ 亮度不变，$A_2$ 变亮}
{闭合 S，$A_1$ 缓慢变亮}
{断开 S 的瞬间，通过 $A_1$ 的电流方向改变}
{断开 S 的瞬间，通过 $A_2$ 的电流方向改变}
%% answer: BD
%% memo:
\memoanswer{
由图可知电路结构为：$L$ 与 $A_1$ 串联为一条支路，$R$ 与 $A_2$ 串联为另一条支路，两条支路并联后，再与电源、开关 S、滑动变阻器 $R_1$ 串联。

A．闭合 S 稳定后缓慢减小 $R$ 的阻值，并联总电阻减小，总电流增大，电源内阻 $r$ 和 $R_1$ 的分压增大，并联部分的电压会减小，因此 $A_1$ 的电流减小、亮度变暗；同时并联电压减小的过程中，线圈 $L$ 的自感会阻碍电流的减小，使 $A_1$ 亮度逐渐变暗；总电流变大，则流过 $A_2$ 支路的电流变大，$A_2$ 变亮，故 A 错误；

B．闭合 S 瞬间，线圈 $L$ 的自感电动势会阻碍自身支路电流的增大，因此 $A_1$ 的电流是从 0 开始逐渐增大的，$A_1$ 缓慢变亮，故 B 正确；

C．闭合 S 稳定时，$L$ 支路的电流方向是从左向右经过 $A_1$；断开 S 瞬间，线圈 $L$ 作为自感电源维持自身电流大小、方向不变，因此通过 $A_1$ 的电流方向依然是从左向右，没有发生改变，故 C 错误；

D．闭合 S 稳定时，$A_2$ 的电流方向是从左向右；断开 S 后，电源被切断，$L$、$A_1$、$A_2$、$R$ 构成闭合的自感回路，回路电流方向为 $L$ 右端 $\to A_1$ 右端 $\to A_2$ 右端 $\to A_2$ 左端 $\to R\to L$ 左端，因此通过 $A_2$ 的电流方向变为从右向左，和原方向相反，发生了改变，故 D 正确。

故选 BD。
}

\item
%% number: 10
%% paperName: 2026年江西高考第 10 题 6 分
%% typeId: 2
%% type: 多选题
%% chapter: 电磁感应
%% point: 电磁感应中的变加速
%% method: 无
%% score: 6
%% degree: 350
%% duplicateId: 0
%% body:
某类电磁装置对物体的作用力源于磁场的空间分布随时间变化，这种变化可等效为磁场区域的空间平移。如图所示，沿水平直线轨道上存在足够多边长为 $10\,\Um$ 的正方形匀强磁场区域，相邻正方形区域磁感应强度大小均为 $5\,\UT$，方向相反，整个磁场区域沿轨道水平向右平移的速度可以调控，其瞬时速度记为 $u$。物体正下方固定连接一水平单匝边长为 $10\,\Um$ 的正方形线圈，电阻为 $1\,\UO$，物体（含线圈）质量为 $1.0\times10^{4}\,\Ukg$，其水平向右运动的瞬时速度记为 $v$，运动过程中线圈完全处于磁场区域。$u$ 与 $v$ 之间的关系为 $u=v+k\sin\left(\frac{\pi}{3}t\right)+b$（$t$ 为时刻，$k$、$b$ 为非负常数），采用国际单位制。忽略线圈自感、空气阻力和一切摩擦。下列选项正确的是\xzanswer{AD}
\onepicture[w=9cm]{figs/10.png}
\fourchoices[answer=AD]
{若 $k<b$，则物体始终做加速直线运动}
{若 $k=1$，$b=0$，则线圈在 $0\sim3\,\Us$ 内产生的焦耳热为 $3.0\times10^{4}\,\UJ$}
{若 $k=0$，$b=1$，则线圈中电流发生周期性变化，其频率为 $0.1\,\UHz$}
{若 $k=1$，$b=1$，则物体在 $0\sim6\,\Us$ 内 $v$ 的增量为 $6\,\Ums$}
%% answer: AD
%% memo:
\memoanswer{
A．若 $k<b$，则 $u-v=k\sin\left(\frac{\pi}{3}t\right)+b\ge-k+b>0$，即始终有 $u>v$。由电磁驱动可知物体始终做加速直线运动，故 A 正确；

B．若 $k=1$，$b=0$，则 $u-v=\sin\left(\frac{\pi}{3}t\right)$。物体相对磁场的切割速度 $v_{\text{相}}=v-u$ 满足正弦函数，故产生正弦式电流，峰值电流
\[
I_m=\frac{2BLv_m}{R}
\]
根据 $u-v=\sin\left(\frac{\pi}{3}t\right)$ 可知相对速度的最大值 $v_m=1\,\Ums$，可得 $I_m=100\,\UA$。在 $0\sim3\,\Us$ 内焦耳热
\[
Q=\left(\frac{I_m}{\sqrt{2}}\right)^2Rt=1.5\times10^{4}\,\UJ
\]
故 B 错误；

C．若 $k=0$，$b=1$，则 $u-v=1\,\Ums$。相对切割速度不变，产生感应电动势大小恒定不变，电流大小恒定不变，方向周期性变化，且周期为
\[
T=\frac{2L}{v_{\text{相}}}=\frac{20}{1}\,\Us=20\,\Us
\]
则频率为
\[
f=\frac{1}{T}=0.05\,\UHz
\]
故 C 错误；

D．若 $k=1$，$b=1$，则 $u-v=\sin\left(\frac{\pi}{3}t\right)+1\ge0$，即物体始终做加速运动。感应电动势大小
\[
E=2BL(u-v)
\]
安培力大小
\[
F_A=2B\cdot\frac{E}{R}\cdot L=\frac{4B^2L^2(u-v)}{R}=10^{4}\left[\sin\left(\frac{\pi}{3}t\right)+1\right]
\]
物体加速度
\[
a=\frac{F_A}{m}=\sin\left(\frac{\pi}{3}t\right)+1
\]
可得 $a-t$ 图像如图所示。
\onepicture[w=6cm]{figs/10_an}
结合 $a-t$ 图像，在 $0\sim6\,\Us$ 内（即一个周期内）与 $t$ 轴包围面积即为 $v$ 的增量，可得 $\Delta v=6\,\Ums$，故 D 正确。

故选 AD。
}

\item
%% number: 11
%% paperName: 2026年江西高考第 11 题 6 分
%% typeId: 4
%% type: 实验题
%% chapter: 电路
%% point: 电表改装
%% method: 无
%% score: 6
%% degree: 650
%% duplicateId: 0
%% body:
某小组利用一个量程为 $100\,\UuA$、内阻为 $1.3\,\UkO$ 的表头（$\UuA$）和两个电阻箱（\includesvg[height=1.1em]{figs/11a}）改装成量程为 $0\sim200\,\UmV$ 和 $0\sim2\,\UV$ 双量程电压表，并进行校准。
\begin{enumerate}
	\item
	改装电压表的设计。请在图（a）虚线框内完成电路设计图，并在其中标出所用电阻的阻值。
	\item
	改装电压表（$2\,\UV$ 量程）的校准。
	\begin{steps}
		\item
		用改装电压表和数字电压表同时测量定值电阻两端电压。现有改装电压表、数字电压表、定值电阻、滑动变阻器、直流电源（$3\,\UV$）、开关和导线，请在图（b）中画线完成电路实物连接。
		\item
		检查电路连接无误后，应先\tkanswer[6cm]{使滑动变阻器滑片移到使定值电阻两端电压为零的位置}，再打开电源，闭合开关，测量数据，进行校准。
	\end{steps}
\end{enumerate}
\twopicture[showlabel=false, align=right, h=4.2cm]
{figs/11b.png}
{figs/11c.png}
%% answer:
\jdanswer{
（1）$R_1=0.7\,\UkO$，$R_2=18\,\UkO$。

（2）①见解析图；②使滑动变阻器滑片移到使定值电阻两端电压为零的位置。
}
%% memo:
\memoanswer{
（1）$200\,\UmV$ 量程的分压电阻
\[
R_1=\frac{U_1}{I_g}-R_g=\frac{0.2\,\UV}{0.1\,\UmA}-1.3\,\UkO=0.7\,\UkO
\]
$2\,\UV$ 量程的分压电阻
\[
R_1+R_2=\frac{U_2}{I_g}-R_g=\frac{2\,\UV}{0.1\,\UmA}-1.3\,\UkO=18.7\,\UkO
\]
可得 $R_2=18\,\UkO$。分压电阻与表头串联，如图所示。
\onepicture[w=6cm]{figs/11_an1}

（2）校准电压表需要让改装电压表和标准数字电压表并联在被测定值电阻两端，同时滑动变阻器接成分压式（因为电源 $3\,\UV$ 大于改装表 $2\,\UV$ 量程，分压式可以实现 $0\sim2\,\UV$ 连续调压，完成全量程校准）：将滑动变阻器的两个下接线柱接电源两端，滑动端和其中一个下接线柱作为输出端，给定值电阻供电；改装电压表选择 $2\,\UV$ 量程接线柱，和数字电压表正负极对应并联在定值电阻两端，如图所示。
\onepicture[w=6cm]{figs/11_an2}

为了保证电路安全，通电前让待测部分电压从 0 开始逐渐增大，避免初始电压过大损坏改装表头，即检查电路连接无误后，应先使滑动变阻器滑片移到使定值电阻两端电压为零的位置，再打开电源，闭合开关，测量数据，进行校准。
}

\item
%% number: 12
%% paperName: 2026年江西高考第 12 题 10 分
%% typeId: 4
%% type: 实验题
%% chapter: 光学
%% point: 光的折射
%% method: 无
%% score: 10
%% degree: 450
%% duplicateId: 0
%% body:
某小组探究盐水浓度对折射率的影响。
\begin{enumerate}
	\item
	首先在纸带上画一条直线，并在直线上作两个明显的标记点 $P_1$、$P_2$，用游标卡尺测出两标记点的距离 $s$，如图（a）所示，读出 $s=$ \tkanswer[1.5cm]{2.325} $\Ucm$。然后将纸带固定在空水槽底部。
	\item
	铁架台和水槽放置在同一水平桌面，激光笔安装于铁架台上。调整实验装置，使激光笔射出的光线与 $P_1P_2$ 连线在同一竖直面内，光线与竖直方向成 $45\Udu$ 角，且光线照射的光斑恰好在标记点 $P_1$，此后装置固定不动，如图（b）所示。向水槽里缓慢加入被测液体，直至光斑沿 $P_1P_2$ 连线移动到标记点 $P_2$，测量并记录液面高度 $h$。如图（c）所示，以 $P_1$ 为坐标原点、沿 $P_1P_2$ 方向为 $x$ 轴建立直角坐标系，虚线为液面所在位置。请在图（c）中画出液面高度为 $h$ 时，激光由空气入射到液体的光路图（入射光线、法线和折射光线）。
	\item
	由折射定律可得，液体的折射率 $n=$ \tkanswer[4cm]{$\frac{\sqrt{2(h-s)^2+2h^2}}{2(h-s)}$}（用 $h$，$s$ 表示）。
	\item
	该小组配制了 5 种不同浓度 $c$ 的盐水，依据上述方案重复实验，所得数据点如图（d）所示。为得到折射率 $n$ 随浓度 $c$ 的变化规律，请在图（d）中作出 $n-c$ 图线。
	\item
	从图（d）中可得盐水折射率 $n$ 与其浓度 $c$ 的函数关系式为 $n=$ \tkanswer[2.7cm]{$0.213c+1.326$}。（式中数字保留三位小数）
\end{enumerate}
\onepicture[align=right, w=6.5cm]{figs/12a.png}
\onepicture[align=right, w=11.5cm]{figs/12b.png}
\onepicture[align=right, w=6cm]{figs/12c.png}
%% answer:
\jdanswer{
\begin{enumerate}
	\item
	$2.325\,\Ucm$
	\item
	光路图见解析。
	\item
	$\frac{\sqrt{2(h-s)^2+2h^2}}{2(h-s)}$
	\item
	$n-c$ 图线见解析。
	\item
	$n=0.213c+1.326$
\end{enumerate}
}
%% memo:
\memoanswer{
（1）20 分度游标尺最小刻度为 $0.05\,\Umm$，故读数
\[
s=23\,\Umm+5\times0.05\,\Umm=2.325\,\Ucm
\]
（2）根据题意可得激光由空气入射到液体的光路图如图所示。
\onepicture[w=5cm]{figs/12_an1}

（3）折射角 $\theta$ 满足
\[
\sin\theta=\frac{h-s}{\sqrt{(h-s)^2+h^2}}
\]
折射率
\[
n=\frac{\sin45^\circ}{\sin\theta}=\frac{\sqrt{2(h-s)^2+2h^2}}{2(h-s)}
\]
（4）根据描点法可得 $n-c$ 图线如图所示。
\onepicture[w=6cm]{figs/12_an2}

（5）由图可知 $n-c$ 图线斜率
\[
k=\frac{1.38-1.33}{0.255-0.02}\approx0.213
\]
截距 $b=1.326$，可得盐水折射率 $n$ 与其浓度 $c$ 的函数关系式为 $n=0.213c+1.326$。
}

\item
%% number: 13
%% paperName: 2026年江西高考第 13 题 8 分
%% typeId: 6
%% type: 计算题
%% chapter: 气体性质
%% point: 热力学第一定律
%% method: 无
%% score: 8
%% degree: 650
%% duplicateId: 0
%% body:
西汉青铜雁鱼灯构思巧妙，体现了我国古代工匠的智慧，如图（a）所示。假设灯内一定质量的理想气体吸收了热量 $Q$，从初态 $A$ 经历压强为 $p_0$ 的等压过程变化到末态 $B$，如图（b）所示。已知气体初态温度为 $T_0$、体积为 $V_0$，末态体积为 $V_1$。求：
\begin{enumerate}
	\item
	气体在末态 $B$ 时的温度 $T_1$；
	\item
	该过程中气体内能的变化量 $\Delta U$。
\end{enumerate}
\onepicture[align=right, w=10cm]{figs/13.png}
%% answer:
\jdanswer{
\begin{enumerate}
	\item
	$T_1=\frac{V_1}{V_0}T_0$
	\item
	$\Delta U=Q-p_0(V_1-V_0)$
\end{enumerate}
}
%% memo:
\memoanswer{
（1）设末态 $B$ 时气体温度为 $T_1$。由盖—吕萨克定律，得
\[
\frac{V_0}{T_0}=\frac{V_1}{T_1}
\]
可得
\[
T_1=\frac{V_1}{V_0}T_0
\]
（2）$AB$ 为等压膨胀，外界对气体做功
\[
W=-p_0(V_1-V_0)
\]
由热力学第一定律 $\Delta U=W+Q$，联立可得
\[
\Delta U=Q-p_0(V_1-V_0)
\]
}

\item
%% number: 14
%% paperName: 2026年江西高考第 14 题 12 分
%% typeId: 6
%% type: 计算题
%% chapter: 运动和力
%% point: 连接体
%% method: 无
%% score: 12
%% degree: 550
%% duplicateId: 0
%% body:
如图（a）所示，质量为 $M=2.0\,\Ukg$ 的长方体箱子在水平拉力 $F$ 作用下沿水平面向右加速运动，箱子与水平面间的动摩擦因数 $\mu=0.60$。箱内后壁上有两物块 A、B，其质量均为 $m=0.5\,\Ukg$，与后壁间的动摩擦因数分别为 $\mu_1=0.25$ 和 $\mu_2=0.50$。$F$ 随时间 $t$ 的变化关系如图（b）所示，当 $F=F_1$ 时，A 恰好静止在后壁上，B 静止在后壁上；当 $F=F_2$ 时，B 恰好静止在后壁上。在运动过程中，A、B 始终没有相碰。设最大静摩擦力等于滑动摩擦力，忽略空气阻力，且 $g$ 取 $10\,\Umsq$。求：
\begin{enumerate}
	\item
	$F_1$ 的大小；
	\item
	$F_2$ 的大小及在 $t_1\sim t_2$ 的时间内摩擦力对物块 A 所做的功，已知 $t_2-t_1=1.0\,\Us$。
\end{enumerate}
\onepicture[align=right, w=10cm]{figs/14.png}
%% answer:
\jdanswer{
\begin{enumerate}
	\item
	$F_1=138\,\UN$
	\item
	$F_2=76.5\,\UN$，$W=-6.25\,\UJ$
\end{enumerate}
}
%% memo:
\memoanswer{
（1）当 $F=F_1$ 时，对 A 有 $\mu_1N_1=mg$，$N_1=ma_1$，得
\[
a_1=40\,\Umsq
\]
对整体有
\[
F_1-\mu(M+2m)g=(M+2m)a_1
\]
得 $F_1=138\,\UN$。

（2）当 $F=F_2$ 时，A 开始向下滑动，对 B 有 $\mu_2N_2=mg$，$N_2=ma_2$，得
\[
a_2=20\,\Umsq
\]
对 A 有 $f_1=\mu_1ma_2$，得 $f_1=2.5\,\UN$。

对箱子有 $N_3=Mg+f_1'+mg$，其中 $f_1'=f_1$，得 $N_3=27.5\,\UN$。

对整体有
\[
F_2-\mu N_3=(M+2m)a_2
\]
得 $F_2=76.5\,\UN$。

对 A 有
\[
mg-\mu_1ma_2=ma_y
\]
得 $a_y=5\,\Umsq$。

有
\[
y=\frac{1}{2}a_y(t_2-t_1)^2,\qquad W=-f_1y
\]
联立可得 $W=-6.25\,\UJ$。
}

\item
%% number: 15
%% paperName: 2026年江西高考第 15 题 18 分
%% typeId: 6
%% type: 计算题
%% chapter: 磁场
%% point: 带电粒子在复合场中的运动
%% method: 无
%% score: 18
%% degree: 450
%% duplicateId: 0
%% body:
冷靶反冲离子动量谱仪可通过测量粒子的动量反推出碰撞动力学过程。其简化装置如图所示，在一个接近真空的圆柱形腔体内以其中心 $O$ 点为原点建立 $Oxyz$ 坐标系，其中 $z$ 轴沿腔体轴线方向。腔内存在沿 $z$ 轴正方向的匀强磁场和沿 $z$ 轴负方向的匀强电场，磁感应强度和电场强度大小分别为 $B$ 和 $E$。腔体顶部（$z=+L$）和底部（$z=-L$）分别装有电子和离子探测器，探测器面积足够大。一束激光与某个超冷中性气体分子在 $O$ 点处发生作用，以类似爆炸的方式产生一个电子和一个反冲离子，电子和反冲离子的质量分别为 $m$、$M$，电荷量分别为 $-e$、$e$。该反冲离子从 $O$ 点运动到离子探测器上的时间为 $\sqrt{\frac{2ML}{eE}}$，打到离子探测器上的位置坐标为 $(x_0,y_0,-L)$。因为 $\frac{M}{m}$ 较大，所以忽略磁场对反冲离子的作用。忽略激光与超冷中性气体分子的动量，不考虑电子和反冲离子所受重力的影响、电子和反冲离子之间的库仑相互作用。（$e$ 为元电荷）
\begin{enumerate}
	\item
	求电子在 $O$ 点处沿 $x$ 轴方向的动量大小；
	\item
	求反冲离子刚打到离子探测器上时的动量大小；
	\item
	若 $y_0=\sqrt{3}x_0$，且 $E$ 与 $B$ 的关系满足 $E=\frac{2eL}{121\pi^2m}B^2$，求电子打到电子探测器上的位置坐标（用 $x_0$、$M$、$m$ 和 $L$ 表示）。
\end{enumerate}
\onepicture[align=right, w=5cm]{figs/15.png}
%% answer:
\jdanswer{
\begin{enumerate}
	\item
	$|x_0|\sqrt{\frac{eEM}{2L}}$
	\item
	$\sqrt{\frac{eEM(x_0^2+y_0^2+4L^2)}{2L}}$
	\item
	$\left(\frac{2\sqrt{3}x_0}{11\pi}\sqrt{\frac{M}{m}},-\frac{2x_0}{11\pi}\sqrt{\frac{M}{m}},L\right)$
\end{enumerate}
}
%% memo:
\memoanswer{
（1）反冲离子只在电场力下做抛体运动，$z$ 轴方向有
\[
-L=v_{z0}t-\frac{1}{2}at^2
\]
其中 $a=\frac{eE}{M}$，$t=\sqrt{\frac{2ML}{eE}}$，可得 $v_{z0}=0$。

即反冲离子做类平抛运动，$x$ 轴方向上有 $x_0=v_xt$，可得
\[
v_x=x_0\sqrt{\frac{eE}{2ML}}
\]
爆炸过程系统动量守恒，故在 $x$ 轴方向上有 $p_{ex}+Mv_x=0$，可得电子在 $O$ 点处沿 $x$ 轴方向的动量大小
\[
|p_{ex}|=|x_0|\sqrt{\frac{eEM}{2L}}
\]
（2）$y$ 轴方向上有 $y_0=v_yt$，可得
\[
v_y=y_0\sqrt{\frac{eE}{2ML}}
\]
$z$ 轴方向上有 $L=\frac{v_z}{2}\cdot t$，可得
\[
v_z=2L\sqrt{\frac{eE}{2ML}}
\]
故末动量大小 $p=M\sqrt{v_x^2+v_y^2+v_z^2}$，可得
\[
p=\sqrt{\frac{eEM(x_0^2+y_0^2+4L^2)}{2L}}
\]
（3）电子沿 $z$ 轴方向做匀加速直线运动，有
\[
L=\frac{1}{2}\cdot\frac{eE}{m}\cdot t'^2
\]
可得
\[
t'=\frac{11\pi m}{eB}
\]
爆炸过程由动量守恒得
\[
u_x=-\frac{M}{m}v_x=-\frac{x_0eB}{11\pi m}\sqrt{\frac{M}{m}},\qquad
u_y=-\frac{M}{m}v_y=\sqrt{3}u_x
\]
电子初速度
\[
u=\sqrt{u_x^2+u_y^2}=\frac{2x_0eB}{11\pi m}\sqrt{\frac{M}{m}}
\]
初速度方向与 $y$ 轴负方向夹角 $\theta$ 满足
\[
\tan\theta=\frac{u_x}{u_y}=\frac{\sqrt{3}}{3}
\]
可得 $\theta=30^\circ$。

在 $Oxy$ 平面内做匀速圆周运动，由 $euB=\frac{mu^2}{r}$，$T=\frac{2\pi r}{u}$，可得
\[
r=\frac{mu}{eB}=\frac{2x_0}{11\pi}\sqrt{\frac{M}{m}},\qquad
T=\frac{2\pi m}{eB}
\]
由于 $t'=\frac{11}{2}T$，故
\[
x=2r\cos\theta=\frac{2\sqrt{3}x_0}{11\pi}\sqrt{\frac{M}{m}},\qquad
y=-2r\sin\theta=-\frac{2x_0}{11\pi}\sqrt{\frac{M}{m}}
\]
\onepicture[w=5cm]{figs/15_an}
故电子打到电子探测器上的位置坐标为
\[
\left(\frac{2\sqrt{3}x_0}{11\pi}\sqrt{\frac{M}{m}},-\frac{2x_0}{11\pi}\sqrt{\frac{M}{m}},L\right)
\]
}

\end{enumerate}

\end{document}
